\documentclass[letterpaper,12pt]{article}
\begin{document}
\newcommand{\TBC}{\framebox{\textbf{TO BE COMPLETED}}}
\newcommand{\be}{\begin{enumerate}}
\newcommand{\ee}{\end{enumerate}}
\newcommand{\bi}{\begin{itemize}}
\newcommand{\ei}{\end{itemize}}
\newcommand{\bd}{\begin{description}}
\newcommand{\ed}{\end{description}}

\newtheorem{definition}{Definition}


\newcommand{\KV}{\mathtt{KV}}

\newcommand{\nHKV}{{n_H}_{\KV}}
\newcommand{\nDKV}{{n_D}_{\KV}}
\newcommand{\tet}{\mathtt{tet}}
\newcommand{\KC}{\mathtt{KC}}

\newcommand{\xb}{\mathtt{xb}}

\newcommand{\wq}{\mathtt{wq}}
\newcommand{\wk}{\mathtt{wk}}
\newcommand{\wv}{\mathtt{wv}}

\newcommand{\wrms}{\mathtt{wrms}}

\newcommand{\rmsnorm}{\mathtt{rmsnorm}}
\newcommand{\vecmatmul}{\mathtt{vecmatmul}}

\section{Conventions and Definitions}
\subsection{Definitions}

\begin{definition}
  \label{scalar}
  A scalar, \(M_0\) is defined using a type \(t_0\)
\end{definition}

\begin{definition}
  \label{1D_matrix}
  A vector or 1D matrix, \(M_1\), is defined using a type \(t_1\) and positive integer
\(n_1\), representing the dimensions/shape.
\begin{itemize}
  \item \(M_1[j]\) is a scalar of type {\tt T} of size \(s_1\)
\item \(M_1\) is stored as a linear array at address \(m_1\)
  such that 
    \begin{enumerate} 
\item Address of \(M_1[j]  = m_1 + s_1 \times j\)
\end{enumerate}
\end{itemize}
\end{definition}

\begin{definition}
  \label{2D_matrix}
  A 2D matrix, \(M_2\),  is defined using a type \(t_2\) and positive integers 
\(n_X, n_Y\), representing the dimensions/shape.
\begin{itemize}
  \item \(M_2[i]\) is a 1D matrix of length \(n_Y\) and type \(t_2\)
  \item \(M_2[i, j]\) is a scalar of type \(t_2\) and size \(s_2\)
  \item \(M_2\) is stored as a linear array at address \(m_2\)
  such that 
    \begin{enumerate} 
      \item Address of \(M_2[i] = m_2 + s_2 \times (i \times n_Y)\)
      \item Address of \(M_2[i, j] = m_2 + 
        s_2 \times (i \times n_Y + j)\)
\end{enumerate}
\end{itemize}
\end{definition}

\begin{definition}
  \label{3D_matrix}
  A 3D matrix, {\tt M}, is defined using a type {\tt T} and positive integers 
\(n_X, n_Y, n_Z\), representing the dimensions/shape.
\begin{itemize}
  \item \(\mathtt{M}[x]\) is a 2D matrix of dimensions \(n_X, n_Y\) and type
    {\tt T}
  \item \(\mathtt{M}[x, y]\) is a vector of length \(n_S\) and type {\tt T}
\item \(\mathtt{M}[x, y, z]\) is a scalar of type {\tt T}
\item {\tt M} is stored as a linear array at address \(m_{\mathtt{M}}\)
  such that 
    \begin{enumerate} 
    \item Address of \(\mathtt{M}[x] = 
  m_{\mathtt{M}} + x \times (n_Y \times n_Z)\)
\item Address of \(\mathtt{M}[x, y]  = 
  m_{\mathtt{M}} + (x \times n_Y \times n_Z) + (y \times n_Z)\)
\item Address of \(\mathtt{M}[x, y, z] = 
  m_{\mathtt{M}} + (x \times n_X \times n_Y) + (y \times n_Z) + z\)
\end{enumerate}
\end{itemize}
\end{definition}



\clearpage
\subsection{Data Structures}
See Section~\ref{state} and ~\ref{weights}.
\subsection{Dimensions}
\label{dimensions}
The constants in Table~\ref{shapes} define the sizes of the various arrays
described later. These are from the {\tt Config struct}.
\begin{table}[hb]
  \begin{tabular}{|l|l|l|} \hline \hline 
    {\bf Math} & {\bf Code} & {\bf Explanation} \\ \hline 
    \(n_L\)  & {\tt n\_layers} & number of layers in the network  \\ \hline 
    \(n_S\)  & {\tt seq\_len} & max sequence length \\ \hline 
    \(n_D\) & {\tt dim} &  transformer dimension \\ \hline 
    \(n_H\) & \verb+n_heads+ & number of query heads \\ \hline
    \(\nDKV\)  & {\tt kv\_dim} & \((n_D \times \nHKV)/n_H\) \\ \hline
     XXXX   & {\tt hidden\_dim} &  for ffn layers \\ \hline 
    \(n_V\) &  {\tt vocab\_size} &  vocabulary size \\ \hline 
    \(\nHKV\)  & \verb+n_kv_heads+ & number of key/value heads \\ \hline
    \(n_{H_0}\) & \verb+n_heads * head_size+ & \TBC \\ \hline 
    \(n_{H_1}\) & \verb+n_kv_heads * head_size+ & \TBC \\ \hline 
    \hline 
\end{tabular}
    \caption{Dimensions}
    \label{shapes}
\end{table}

\section{State}
\label{state}
The state is represented by the data structures listed in
Table~\ref{tbl_state_data}
\begin{table}[h]
  \begin{tabular}{|l|l|l|l|l|l|} \hline \hline
    {\bf Name} & {\bf Shape} & {\bf Comment} & \(n_X\) & \(n_Y\) & \(n_Z\) \\ \hline
    {\tt x} & 1D array & activation at current timestamp &  \(n_D\) & & \\ \hline
    {\tt xb} & 1D array & like {\tt x} but inside residual branch & \(n_D\) & &
    \\ \hline
    {\tt xb2} & 1D array & like {\tt xb} & \(n_D\) & & \\ \hline
    \hline
    {\tt sq} & 1D array & \(n_H\) & & \\ \hline
    {\tt sk} & 1D array & \(n_D\) & & \\ \hline
    {\tt sk} & 1D array & \(n_D\) & & \\ \hline
\hline
    {\tt KC} & 3D array & keys of kv-cache & \(n_L\) & \(n_S\) & \(n_S\)  \\
    \hline 
    {\tt KV} & 3D array & values of kv-cache & \(n_L\) & \(n_S\) & \(n_S\)  \\ \hline 
    \hline
\end{tabular}
  \caption{Data structures for state}
  \label{tbl_state_data}
\end{table}

\section{Weights}
\label{weights}
The weights are represented by the data structures listed in
Table~\ref{tbl_weights_data}

\begin{table}
  \begin{tabular}{|l|l|l|l|l|l|} \hline \hline
    {\bf Name} & {\bf Shape} & {\bf Comment} & \(n_X\) & \(n_Y\) & \(n_Z\) \\ \hline
    {\tt tet} & 2D array & Lookup table maps token IDs to vectors  & \(n_V\) & \(n_D\) & \\ \hline 
    {\tt wk} & 3D array & key projections & \(n_L\) & \(n_D\) & \(n_{H_1}\) \\ \hline 
    {\tt wv} & 3D array & value projections & \(n_L\) & \(n_D\) & \(n_{H_1}\) \\ \hline 

    \hline
\end{tabular}
  \caption{Data structures for weights}
  \label{tbl_weights_data}
\end{table}

  
\section{Code}

Given 
\begin{enumerate}
\item an integer \(p\) where \(0 \leq p < n_P\)
\item an integer \(t\) where \(0 \leq t < n_V\)
  \end{enumerate}

%%--------------------------------------------
  \subsection{rmsnorm}
\begin{figure}[h]
  \centering
\fbox{
  \begin{minipage}{6.5in}
  \begin{tabbing} \hspace*{0.25in} \= \hspace*{0.25in} \=  \kill
    {\tt function} rmsnorm(\+ \\
    \(x\), /* type is (\(M_1, n\), {\tt float}) */ \\ 
    \(w\)  /* type is (\(M_1, n\), {\tt float}) */ \- \\
    ) \\
    {\bf begin} \+ \\
    \(y\) /* type is (\(M_1, n\), {\tt float}) */ \\
    \(\epsilon = 10^{-5}\) /* type is (\(M_0\), {\tt float}) */ \\
    \(s \leftarrow \frac{1}{\sqrt{\epsilon + \frac{\sum {x_i}^2}{n}}}\) \\
    \(y_i \leftarrow w_i \times x_i \times s\) \\ 
    {\bf return} y \- \\ 

    {\bf end} \\ 
\end{tabbing}
  \caption{Pseudo-code for \(\rmsnorm()\) function}
  \label{fn_rmsnorm}
  \end{minipage}
}
\end{figure}
\clearpage
%%--------------------------------------------
  \subsection{vecmatmul}
\begin{figure}[h]
  \centering
\fbox{
  \begin{minipage}{6.5in}
  \begin{tabbing} \hspace*{0.25in} \= \hspace*{0.25in} \=  \kill
    {\tt function} vecmatmul(\+ \\
    \(x\), /* type is (\(M_1, n\), {\tt float}) */ \\ 
    \(w\),  /* type is (\(M_2, m, n\), {\tt float}) */ \\
    \(n\), /* type is (\(M_0\), {\tt float}) */ \\
    \(m\) /* type is (\(M_0\), {\tt float}) */ \- \\
    ) \\
    {\bf begin} \+ \\
    \(y\) /* type is (\(M_1, m\), {\tt float}) */ \\
    \(y_i \leftarrow \sum_{i=0}^{i=n-1} w[i][j] \times x[j]\) \\ 
    {\bf return} y \- \\ 
    {\bf end} \\ 
\end{tabbing}
  \caption{Pseudo-code for {\tt vecmatmul()} function}
  \label{fn_vecmatmul}
  \end{minipage}
}
\end{figure}
\clearpage
%%--------------------------------------------
  \subsection{rope}
\begin{figure}[h]
  \centering
\fbox{
  \begin{minipage}{6.5in}
  \begin{tabbing} \hspace*{0.25in} \= \hspace*{0.25in} \=  \kill
    {\tt function} rope(\+ \\
    \(q\), /* type is (\(M_1, n\), {\tt float}) */ \\ 
    \(k\),  /* type is (\(M_1, m, n\), {\tt float}) */ \\
    \(n\), /* type is (\(M_0\), {\tt float}) */ \\
    \(m\) /* type is (\(M_0\), {\tt float}) */ \- \\
    ) \\
    {\bf begin} \+ \\
    \(y\) /* type is (\(M_1, m\), {\tt float}) */ \\
    \(y_i \leftarrow \sum_{i=0}^{i=n-1} w[i][j] \times x[j]\) \\ 
    {\bf return} y \- \\ 
    {\bf end} \\ 
\end{tabbing}
  \caption{Pseudo-code for {\tt rope()} function}
  \label{fn_rope}
  \end{minipage}
}
\end{figure}
\clearpage
%%--------------------------------------------
\subsection{forward}
\begin{figure}[h]
  \centering
\fbox{
  \begin{minipage}{6.5in}
  \begin{tabbing} 
    \hspace*{0.25in} \= \hspace*{0.25in} \=  
    \hspace*{0.25in} \= \hspace*{0.25in} \=  
    \kill
    {\bf function} {\tt foo} ( \+ \\
    {\bf integer} \(t\), \\
    {\bf integer} \(p\) \- \\
    ) \\
    {\bf begin} \+ \\
{\bf foreach} \(l \in \{0, 1, \ldots n_L-1\}\) {\bf do}  \+ \\ 

  %% copy the token embedding into x
  %% float* tet_ptr = mcr_2d_to_1d(w->token_embedding_table, token, ispc_dim);
  %% memcpy(x, tet_ptr, ((size_t)dim*sizeof(float)));
    \(x \leftarrow \tet[t]\) \\  
    % rmsnorm(s->xb, x, w_rms_att_l, dim);
    \(\xb \leftarrow \rmsnorm(x, \wrms[l])\) \\ 
    \(\mathtt{q}[l] \leftarrow \vecmatmul(\xb, \wq[l])\) \\ 
    \(\mathtt{k}[l] \leftarrow \vecmatmul(\xb, \wk[l])\) \\ 
    \(\mathtt{v}[l] \leftarrow \vecmatmul(\xb, \wv[l])\) \- \\ 
%%        matmul(s->q, s->xb, w->wq + l*dim*dim, dim, dim);
%%        matmul(s->k, s->xb, w->wk + l*dim*kv_dim, dim, kv_dim);
%%        matmul(s->v, s->xb, w->wv + l*dim*kv_dim, dim, kv_dim);
    \\ 
    {\bf endfor}  \- \\
    {\bf end} 
\end{tabbing}
  \end{minipage}
}
  \caption{Pseudo-code for {\tt forward()} function}
  \label{fn_forward}
\end{figure}
%      matmul(key_ptr, s->xb, w_k, dim, kv_dim);
%      matmul(val_ptr, s->xb, w_v, dim, kv_dim);
\clearpage

\section{Instructions to LLM}
For each of the functions listed in Table~\ref{list_of_fns}
\begin{enumerate}
  \item Write an optimized C function using AVX2 intrinsics and SIMD
    instructions whenever possible.
  \item Return a {\tt .c} file --- the definition
  \item Return a {\tt .h} file --- the declaration
\end{enumerate}
\begin{table}
\begin{tabular}{|l|l|} \hline \hline
  {\bf Function} & Section Implemented  \\ \hline \hline
  {\tt rmsnorm} & Section~\ref{fn_rmsnorm} \\ \hline
  {\tt vecmatmul} & Section~\ref{fn_vecmatmul} \\ \hline
  {\tt rope} & Section~\ref{fn_rope} \\ \hline
%%  {\tt forward} & Section~\ref{fn_forward} \\ \hline
\end{tabular}
  \caption{Listing of functions}
  \label{list_of_fns}
\end{table}

\end{document}
